\documentclass[11pt]{article}
\usepackage[margin=1in]{geometry}
\usepackage[T1]{fontenc}
\usepackage[utf8]{inputenc}
\usepackage{lmodern}
\usepackage{amsmath,amssymb,amsthm,mathtools}
\usepackage{booktabs,array,longtable}
\usepackage{enumitem}
\usepackage{setspace}
\usepackage{microtype}
\usepackage{csquotes}
\usepackage{hyperref}
\usepackage[backend=biber,style=authoryear,natbib=true,maxcitenames=2,maxbibnames=99]{biblatex}
\usepackage{titlesec}
\hypersetup{colorlinks=true,linkcolor=blue,citecolor=blue,urlcolor=blue}
\setstretch{1.08}
\newtheorem{proposition}{Proposition}
\newtheorem{lemma}{Lemma}
\newtheorem{assumption}{Assumption}
\newcommand{\E}{\mathbb{E}}
\newcommand{\R}{\mathbb{R}}
\newcommand{\Prob}{\mathbb{P}}
\newcommand{\one}{\mathbf{1}}
\DeclareMathOperator*{\argmax}{arg\,max}
\titleformat{\section}{\large\bfseries}{\thesection}{0.7em}{}
\titleformat{\subsection}{\normalsize\bfseries}{\thesubsection}{0.7em}{}
\titleformat{\subsubsection}{\normalsize\itshape}{\thesubsubsection}{0.7em}{}

\addbibresource{Ambiguity_Rollover_Crises.bib}

\title{Ambiguity and Sovereign Rollover Crises: A Top-level PhD Research Proposal}
\author{Prepared as an independent proposal draft}
\date{July 2026}

\begin{document}
\maketitle
\begin{abstract}
This standalone proposal (Original Topic 5) states a focused research question, positions it against verified literature, develops a tractable model, derives testable implications, and specifies an empirical design to validate or reject the mechanism. It is self-contained: the preamble is embedded in this file and the references are contained in the local BibTeX file only.
\end{abstract}

\paragraph{Source-integrity note.} References are restricted to publications, working papers, datasets, and institutional sources checked during preparation. Unverified claims are not used as established facts.

\section{Ambiguity, Investor Base Composition, and Sovereign Rollover Crises}

\paragraph{One-sentence pitch.} Rollover crises depend not only on fundamentals and information precision, but also on who holds the debt. This project studies how ambiguity-averse, price-sensitive investors shift the crisis threshold in a sovereign global game.

\subsection*{1. Research question and reviewer positioning}
The core question is:
\begin{quote}
How does ambiguity about public information quality, combined with changes in the investor base, affect sovereign rollover-crisis thresholds and the value of official information disclosure?
\end{quote}
The theoretical foundation is global games. \citet{MorrisShin2003} show how noisy private information can select a unique equilibrium in coordination games. \citet{MorrisShin2004} apply coordination risk to debt pricing. \citet{HeXiong2012} analyze dynamic debt runs with staggered creditors. \citet{Ui2025} adds strategic ambiguity and shows that ambiguous information can have different effects from low-quality information; in debt rollover applications, more ambiguous-quality news can trigger rollover crises because not rolling over is the safe action.

The proposed contribution is to bring this logic into a sovereign-debt environment with a changing investor base. The paper is not simply ``global games plus sovereign debt.'' It adds a measurable state variable: the share of debt held by investors who are more ambiguity-averse, more leveraged, or more sensitive to mark-to-market losses. The comparative statics then connect theory to data.

\subsection*{2. Contribution relative to the literature}
The paper contributes in three ways. First, it develops a sovereign rollover game with heterogeneous ambiguity sets. Second, it maps investor-base composition into the aggregate rollover threshold. Third, it evaluates policy tools - IMF signals, fiscal transparency, and collective action clauses - through their effect on ambiguity rather than only through fundamentals.

A top referee will ask whether ambiguity is observable. The answer is that the model does not require directly observing preferences. It requires proxies for ambiguity sensitivity and belief dispersion: CDS-implied volatility smiles, survey disagreement, bid-ask spreads, investor-base shares, mutual-fund flow sensitivity, and official-sector ownership. These objects do not perfectly identify ambiguity, but they allow disciplined tests of the model's comparative statics.

\subsection*{3. Static rollover game}
A sovereign has short-term debt $B$ due today. There is a continuum of creditors of measure one. Each creditor chooses $a_i\in\{0,1\}$, where $a_i=1$ means roll over and $a_i=0$ means run. The sovereign survives if the mass of runners is below a threshold depending on fundamental $\theta$:
\begin{equation}
\int_0^1 (1-a_i)di \leq \Lambda(\theta), \qquad \Lambda'(\theta)>0. \label{eq:survival}
\end{equation}
Higher $\theta$ means stronger fundamentals and a larger tolerable run.

Creditor $i$ observes
\begin{equation}
x_i=\theta+\sigma_k \varepsilon_i,
\end{equation}
where $k$ indexes investor type. Type $k$ has ambiguity set $\mathcal P_k(x_i)$ over the posterior distribution of $\theta$ and others' actions. The payoff from rolling over is
\begin{equation}
u^R_i=\begin{cases}
R_k-1, & \text{if the sovereign survives},\\
\mathcal R_k-1, & \text{if the sovereign fails},
\end{cases}
\end{equation}
with $R_k>1>\mathcal R_k$. The payoff from running is normalized to $0$. Under max-min preferences, type $k$ rolls over if
\begin{equation}
\min_{P\in\mathcal P_k(x_i)} E_P[u^R_i]\geq 0. \label{eq:maxmin}
\end{equation}
Let $P_k^{\max}(C\mid x_i)$ be the worst-case posterior probability of crisis. Condition \eqref{eq:maxmin} becomes
\begin{equation}
(1-P_k^{\max}(C\mid x_i))(R_k-1)+P_k^{\max}(C\mid x_i)(\mathcal R_k-1)\geq 0. \label{eq:rollover-condition}
\end{equation}
Equivalently, type $k$ rolls over when
\begin{equation}
P_k^{\max}(C\mid x_i)\leq \bar P_k\equiv \frac{R_k-1}{R_k-\mathcal R_k}. \label{eq:thresholdprob}
\end{equation}
Because $P_k^{\max}(C\mid x_i)$ is decreasing in $x_i$, there exists a cutoff $x_k^*$ such that type $k$ runs iff $x_i<x_k^*$.

\subsection*{4. Investor base and equilibrium threshold}
Let $\mu_{k,t}$ be the share of sovereign debt held by type $k$ investors. Given cutoffs $x_k^*$, the expected mass of runners at fundamental $\theta$ is
\begin{equation}
Run(\theta;\mu)=\sum_k \mu_k \Phi\left(\frac{x_k^*-\theta}{\sigma_k}\right). \label{eq:aggregate-run}
\end{equation}
The crisis threshold $\theta^*$ solves
\begin{equation}
Run(\theta^*;\mu)=\Lambda(\theta^*). \label{eq:crisis-threshold}
\end{equation}

\begin{proposition}[Ambiguity raises the run cutoff]
If a larger ambiguity set raises the worst-case crisis probability $P_k^{\max}(C\mid x)$ for every signal $x$, then the rollover cutoff $x_k^*$ rises. Ambiguity-averse creditors run at stronger private signals.
\end{proposition}

\begin{proposition}[Investor-base fragility]
If type $A$ has a higher run cutoff than type $B$, then increasing $\mu_A$ raises the aggregate crisis threshold $\theta^*$ whenever the left-hand side of \eqref{eq:crisis-threshold} crosses the survival capacity from above.
\end{proposition}

\begin{proposition}[Non-monotone value of public information]
A public signal that reduces ambiguity can lower crisis risk, but a public signal that is precise about bad fundamentals can raise crisis risk. Therefore the welfare effect of transparency depends on whether the policy primarily reduces ambiguity or reveals adverse fundamentals.
\end{proposition}

The third proposition is important. It prevents a naive policy conclusion that all disclosure is always good.

\subsection*{5. Dynamic extension}
The dynamic model adds debt maturity and investor-base evolution:
\begin{equation}
B_{t+1}=\mathcal B(B_t,r_t,PrimaryBalance_t), \qquad \mu_{t+1}=\mathcal M(\mu_t,QT_t,Flows_t,Returns_t). \label{eq:dynamic-investor-base}
\end{equation}
The sovereign spread satisfies
\begin{equation}
spread_t=\Pi(\theta_t,B_t,\mu_t,Ambiguity_t,GlobalRisk_t). \label{eq:spread-ambiguity}
\end{equation}
This dynamic block is deliberately parsimonious. The goal is not a full DSGE model; it is a tractable map from investor composition to rollover fragility.

\subsection*{6. Empirical design}
\paragraph{Data.} The empirical work uses sovereign bond and CDS spreads, CDS option-implied volatility when available, sovereign auction data, investor-base data from BIS/IMF/national debt-management offices, official-sector purchase and QT episodes, IMF Article IV and program dates, ratings announcements, and macro controls.

\paragraph{Reduced-form test 1: investor base and spread sensitivity.}
\begin{equation}
\Delta spread_{i,t}=\alpha_i+\tau_t+\beta\,AmbiguityProxy_{i,t}\times FragileInvestorShare_{i,t}+\Gamma'X_{i,t}+\varepsilon_{i,t}. \label{eq:ambig-panel}
\end{equation}
The model predicts $\beta>0$: ambiguity matters more when the marginal investor is more likely to run.

\paragraph{Reduced-form test 2: public signals.} Around IMF Article IV releases, IMF program news, or major fiscal announcements, estimate
\begin{equation}
\Delta spread_{i,t+h}=\alpha_i+\tau_t+\sum_h \beta_h Signal_{i,t}\times AmbiguityProxy_{i,t-1}+\Gamma'X_{i,t-1}+u_{i,t+h}. \label{eq:signal-lp}
\end{equation}
If the signal mainly reduces ambiguity, spreads should fall more for high-ambiguity countries. If the signal reveals bad fundamentals, the sign may reverse. The paper must classify signals carefully rather than pool all announcements.

\paragraph{Structural validation.} Estimate the model by matching: (i) spread jumps around public signals; (ii) changes in rollover auction tails; (iii) the relationship between investor-base composition and spread volatility; and (iv) crisis episodes where rollover pressure materializes. The key mapping is from observed $\mu_k$ and ambiguity proxies to the implied crisis threshold $\theta^*$.

\subsection*{7. Falsifiable predictions}
\begin{enumerate}[label=\textbf{P\arabic*.},leftmargin=*]
\item Spreads react more to bad public news when fragile investor share is high.
\item Ambiguity-reducing official communication lowers spreads most when prior ambiguity is high and fundamentals are not too weak.
\item Quantitative tightening or official-sector exit raises rollover fragility when it shifts debt toward more price-sensitive investors.
\item Collective action clauses reduce crisis probability more in high-ambiguity environments because they change the payoff from running.
\item CDS-option smiles and forecast dispersion should predict future spread jumps even after controlling for debt, growth, ratings, and VIX.
\end{enumerate}

\subsection*{8. Dissertation structure}
\begin{enumerate}[leftmargin=*]
\item \textbf{Chapter 1: Theory.} Solve the heterogeneous-ambiguity global game and characterize investor-base comparative statics.
\item \textbf{Chapter 2: Empirics.} Estimate how investor-base composition changes the sensitivity of sovereign spreads to ambiguity and public signals.
\item \textbf{Chapter 3: Policy.} Evaluate IMF communication, fiscal transparency, and CAC design as ambiguity-management tools.
\end{enumerate}

\subsection*{9. Reviewer stress test}
The first risk is that ambiguity is not cleanly identified. The paper should be honest: it tests comparative statics using proxies, not revealed preference at the individual-investor level. The second risk is theoretical novelty. To be more than an application of \citet{Ui2025}, the paper must make investor-base composition central and empirically disciplined. The third risk is data availability for creditor composition. The best starting sample is sovereigns with detailed debt-management-office data and liquid CDS markets, not the widest possible emerging-market panel.

\subsection*{10. Bottom line}
This is the most theoretically elegant proposal. It is high-risk because the proofs must be clean and the ambiguity proxies must be credible. If successful, it offers a rigorous bridge between global games, sovereign debt markets, and current policy concerns about the changing marginal holder of government debt.


\newpage
\printbibliography
\end{document}
